% 来源：OEBPS/Text/01_30.xhtml（OPF spine 顺序；仅第一册）
\cleardoublepage
\BookEntry[格里高利·曼昆《经济学原理》（第7版）]{格里高利·曼昆《经济学原理》（第7版）}

\BookCover{images/书单16经济学原理.jpg}{[美] 格里高利·曼昆著| 梁小民、梁砾译| 北京大学出版社2015年版}

\begin{BookSummaryBox}
\textbf{世界上最受欢迎的经济学教材}
\par
\end{BookSummaryBox}

\subsubsection*{这本书谈什么？}

哈佛大学经济学教授曼昆的《经济学原理》(第7版)是目前国内市场上最受欢迎的经济学教材，分“微观经济学”和“宏观经济学”两册。它的最大特点是“学生导向”，与同类教科书相比，本书更强调经济学原理的应用和政策分析，而不是正式的经济学模型。
\par

作者在大部分章节里都提供了案例，以说明经济学原理如何应用于现实经济学问题的分析。为了使读者懂得什么是生活中的经济学，作者在书中还提供了大量的“新闻摘录”。该书已被翻译成20多种语言。
\par

\subsubsection*{芒格为什么谈起这本书？}

2003年10月3日，芒格在“论学院派经济学”的演讲中，谈到学院派经济学的一些明显优点。例如，它生逢其时、生逢其地，自然得到了发展。所有计划经济都崩溃以后，那些自由市场经济或者半自由市场经济都蓬勃发展，这增加了经济学的名声。而关于经济学如何通过汲取其他学科的养分，进行跨学科研究，芒格说道：
\par

\begin{quotation}
在格里高利·曼昆撰写的教材中，我们可以发现，这位经济学家从其他学科吸取所需养分的本领已经非常高明了。我肯定是美国少数在那本书刚出版时就买下来的商人之一，因为那本教材得了一大笔预付稿费。我想弄清楚那个家伙到底做了些什么，怎么能够取得这么大的一笔稿费。所以我就这么凑巧把曼昆这本为大一学生写的教材给看完了。书中列举了许多经济学原理：机会成本是一种超级力量，所有希望获得正确答案的人都可以使用。还有，激励机制也是一种超级力量。
\par

最后还有“公地悲剧”的原理，这个原理是由我的老朋友，加州大学圣塔巴巴拉分校教授加勒特·哈丁提出的。哈丁为经济学引入了一只邪恶的无形之脚，它足以和斯密那只做好事的无形之手相提并论。我认为哈丁的理论使经济学变得更加完善。哈丁当年向我介绍他的理论时，我就知道他这个“公地悲剧”理论迟早会被写进教科书。你们看啊，二十年过去了，它终于被写进了经济学教材。曼昆这种借鉴其他学科、吸取哈丁的理论和其他有用知识的做法是很正确的。
\par
\end{quotation}

\subsubsection*{《经济学原理》作者格里高利·曼昆如是说}

“一个家庭面临这许多决策，它必须考虑到每个成员的能力、努力和愿望，以在其各个成员中配置稀缺资源。”
\par

“由于资源是稀缺的，社会资源的管理就显得尤为重要。”
\par

“社会上的每个人不能达到他希望的最高生活水平。”
\par

“某种东西的成本是为了得到它所放弃的东西。”
\par

“贸易可以使每个人的状况都变得更好。”
\par

“征服政策者绝不能忘记激励，因为许多政策改变了人们面临的成本和收益。”
\par
